\chapter{संरचना निर्धारण: \texorpdfstring{$^{13}$C}{13C} NMR, MS और IR एक साथ}\label{ch:b2:structure-determination}

एक सुगंध प्रयोगशाला को चमेली जैसी गंध वाले एक रंगहीन द्रव की शीशी मिलती है। किसी ने कोई
लेबल नहीं लिखा। \omterm{def:b2:mass-spec-atomic:spectrum}{द्रव्यमान स्पेक्ट्रम} उसका आण्विक सूत्र देता है, अवरक्त स्पेक्ट्रम उसके क्रियात्मक
समूह, कार्बन-13 स्पेक्ट्रम उसका ढाँचा और प्रोटॉन स्पेक्ट्रम यह कि टुकड़े कैसे जुड़े हैं; एक
घंटे के भीतर द्रव का नाम मिल जाता है। यह अध्याय प्रथम वर्ष के खंड की प्रोटॉन NMR और अवरक्त स्पेक्ट्रमिकी
तथा \cref{ch:b2:mass-spec-atomic} की \omterm{def:b2:mass-spec-atomic:spectrum}{द्रव्यमान स्पेक्ट्रममिति} में \omterm{def:b2:structure-determination:c13}{कार्बन-13 NMR} जोड़ता है, और चारों को
एक विधि में बदल देता है।

\begin{recall}
प्रथम वर्ष का खंड: रासायनिक विस्थापन और परिरक्षण, तुल्य प्रोटॉन, समाकलन, चक्रण--चक्रण युग्मन
और $n + 1$ नियम, क्रियात्मक समूहों की अवरक्त तरंग संख्याएँ और अंगुलिछाप क्षेत्र, असंतृप्ति
की मात्रा, तीन स्पेक्ट्रमों से संरचना निकालना। \Cref{ch:b2:mass-spec-atomic}: \omterm{def:b2:mass-spec-atomic:molecular-ion}{आण्विक
आयन}, \omterm{def:b2:mass-spec-atomic:isotope-pattern}{समस्थानिक पैटर्न}, विखंडन।
\end{recall}

\section{कार्बन-13 NMR}

कार्बन-12, मुख्य समस्थानिक, का कोई नाभिकीय चक्रण नहीं है और वह कोई NMR संकेत नहीं देता। कार्बन-13 का चक्रण
प्रोटॉन की तरह $\tfrac12$ है, परंतु वह प्राकृतिक कार्बन का केवल 1.1~\% है और उसका संकेत,
नाभिक-प्रति-नाभिक, प्रोटॉन के संकेत से दुर्बल है; अतः कार्बन स्पेक्ट्रम के लिए अधिक नमूना या अधिक क्रमवीक्षण चाहिए।
% ledger: iso:C13

\begin{definition}[कार्बन-13 NMR]\label{def:b2:structure-determination:c13}
\emph{कार्बन-13 NMR}\index{कार्बन-13 NMR} किसी नमूने के \ce{^{13}C} नाभिकों के अनुनाद अभिलिखित करता है।
इसे प्रायः \emph{विस्तृत-पट्ट वियुग्मन}\index{विस्तृत-पट्ट वियुग्मन} के साथ चलाया जाता है: मापन के दौरान
प्रोटॉनों को उनकी आवृत्तियों के पूरे परास पर विकिरित किया जाता है, जो हर
कार्बन--प्रोटॉन युग्मन को हटा देता है, जिससे हर प्रकार का कार्बन एक अकेली रेखा देता है। \emph{तुल्य
कार्बन}\index{तुल्य कार्बन}, जो अणु की किसी सममिति (या किसी तेज़ घूर्णन) द्वारा आपस में बदलते हैं,
एक ही रेखा देते हैं।
\end{definition}

\begin{proposition}[कोई कार्बन--कार्बन युग्मन नहीं]\label{prop:b2:structure-determination:no-cc-coupling}
प्राकृतिक-प्रचुरता वाले \ce{^{13}C} स्पेक्ट्रम में पड़ोसी कार्बन परमाणुओं के बीच युग्मन
दिखाई नहीं देते।
\end{proposition}

\begin{proof}
दो कार्बनों के बीच युग्मन केवल उन अणुओं में दिखता है जिनमें दोनों \ce{^{13}C} हों। दो दी गई
पड़ोसी स्थितियों के लिए प्रायिकता $a_{13}^2 = 0.011^2 \approx 1.2 \times 10^{-4}$ है: लगभग आठ हज़ार में एक
अणु। हर कार्बन का संकेत लगभग पूरा उन अणुओं से आता है जिनमें उसके
पड़ोसी \ce{^{12}C} हैं, जो स्पेक्ट्रममापी के लिए अदृश्य हैं; युग्मित उपग्रह किसी सामान्य स्पेक्ट्रम के रव से सौ गुना
दुर्बल होते हैं।
\end{proof}

\begin{proposition}[रेखाओं की संख्या]\label{prop:b2:structure-determination:signals}
वियुग्मित \ce{^{13}C} स्पेक्ट्रम \omterm{def:b2:structure-determination:c13}{तुल्य कार्बनों} के हर समुच्चय के लिए एक रेखा दिखाता है, जब तक कि संयोग से दो समुच्चयों का
विस्थापन समान न हो।
\end{proposition}

\begin{proof}[तर्क]
\omterm{def:b2:structure-determination:c13}{तुल्य कार्बन} समान परिवेश में होते हैं, अतः उनका विस्थापन समान होता है; अतुल्य कार्बन
सामान्यतः भिन्न होते हैं। प्रोटॉनों से युग्मन हटा दिए जाने और कार्बनों के बीच युग्मन अनुपस्थित होने पर
(\cref{prop:b2:structure-determination:no-cc-coupling}) कोई चीज़ रेखा को विभाजित नहीं करती। संयोगवश अतिव्यापन
तब होता है जब दो भिन्न कार्बनों के विस्थापन विभेदन से अधिक निकट हों, जो
\qty{128}{ppm} के निकट बेन्ज़ीन वलय के कार्बनों के लिए सामान्य है।
\end{proof}

विस्थापन लगभग \qty{220}{ppm} पर फैले होते हैं, प्रोटॉनों के परास के बीस गुना, अतः अतिव्यापन
प्रोटॉन स्पेक्ट्रमों से दुर्लभ हैं, और केवल विस्थापन ही कार्बन का प्रकार बता देता है। नीचे का चार्ट इस अध्याय के यौगिकों के लिए
मापे गए विस्थापन दिखाता है।

\begin{omfigure}
\begin{tikzpicture}
\begin{axis}[width=0.86\linewidth, height=4.6cm, xmin=0, xmax=220, ymin=-0.6, ymax=4.6, x dir=reverse,
  xlabel={$\delta$ (ppm)}, label style={font=\small}, tick label style={font=\scriptsize},
  ytick={0,1,2,3,4}, yticklabels={ऐल्किल C, C--O, ऐरोमैटिक C, एस्टर C=O, कीटोन C=O},
  xtick={0,20,40,60,80,100,120,140,160,180,200,220}, grid=major, grid style={black!10}]
\addplot[only marks, mark=|, mark size=5pt, thick, omDef] table[x=shift, y=cls] {figdata/out/bachelor-2-nmr-spectra-d.dat};
\end{axis}
\end{tikzpicture}

{\small ब्यूटेन-2-ओन, पेंटेन-2-ओन, एथिल बेन्ज़ोएट और बेन्ज़िल एथेनोएट के मापे गए \ce{^{13}C} विस्थापन,
कार्बन के प्रकार के अनुसार छाँटे गए (CDCl$_3$; पबकेम आँकड़े)। कीटोन कार्बोनिल 209 के निकट होते हैं, एस्टर कार्बोनिल
167--171 के निकट, \omterm{def:b2:huckel:aromatic}{ऐरोमैटिक} कार्बन 128 और 137 के बीच, एस्टर ऑक्सीजन से आबंधित कार्बन 61--66 के निकट, और
ऐल्किल कार्बन 46 से नीचे।}
% figdata: bachelor-2/nmr-spectra.py part d
% ledger: c13:butanone, c13:pentan-2-one, c13:ethylbenzoate, c13:benzylacetate
\end{omfigure}

\begin{proposition}[ऊँचाइयाँ गिनतियाँ नहीं हैं]\label{prop:b2:structure-determination:intensities}
किसी सामान्य वियुग्मित \ce{^{13}C} स्पेक्ट्रम में रेखाओं की ऊँचाइयाँ उनके द्वारा दर्शाए गए
कार्बनों की संख्याओं के समानुपाती नहीं होतीं; हाइड्रोजन-रहित कार्बन दुर्बल रेखाएँ देते हैं।
\end{proposition}

\begin{proof}[तर्क]
सामान्य स्पेक्ट्रम मापन को उससे तेज़ी से दोहराते हैं जितनी तेज़ी से सबसे धीमे कार्बन साम्य में लौटते हैं, और
जुड़े प्रोटॉन से रहित कार्बन धीरे-धीरे लौटते हैं, अतः उनका संकेत आंशिक रूप से संतृप्त हो जाता है। वियुग्मन
प्रोटॉन धारण करने वाले कार्बनों के संकेतों को भी बढ़ाता है, इतनी मात्रा में जो उनके परिवेश पर निर्भर करती है।
ब्यूटेन-2-ओन में कार्बोनिल रेखा चारों में सबसे दुर्बल है, यद्यपि वह भी अन्य रेखाओं की तरह एक ही कार्बन को दर्शाती है।
अतः समाकलन, जो प्रोटॉनों के लिए इतना उपयोगी है, सामान्य कार्बन स्पेक्ट्रमों के लिए प्रयुक्त नहीं होता।
\end{proof}

\section{DEPT}

\begin{definition}[DEPT]\label{def:b2:structure-determination:dept}
\emph{DEPT}\index{DEPT} (ध्रुवण स्थानांतरण द्वारा अविकृत संवर्धन) कार्बन-13 के प्रयोगों का एक समुच्चय है
जो जुड़े प्रोटॉनों वाले कार्बनों के संकेतों को उनकी हाइड्रोजन संख्या के अनुसार एक चिह्न देता है या उन्हें
अनुपस्थित कर देता है। \emph{चतुष्क कार्बन}\index{चतुष्क कार्बन} पर कोई हाइड्रोजन नहीं होता
(चार अन्य परमाणुओं से आबंधित, या H-रहित कोई कार्बोनिल या \omterm{def:b2:huckel:aromatic}{ऐरोमैटिक} कार्बन); वह कोई DEPT संकेत नहीं देता।
\end{definition}

\begin{proposition}[DEPT पढ़ना]\label{prop:b2:structure-determination:dept}
DEPT-135 में \ce{CH} और \ce{CH3} कार्बन ऊपर इंगित करते हैं, \ce{CH2} कार्बन नीचे, और \omterm{def:b2:structure-determination:dept}{चतुष्क कार्बन}
अनुपस्थित होते हैं। DEPT-90 में केवल \ce{CH} कार्बन दिखते हैं।
\end{proposition}

\admitted

प्रयोग प्रोटॉनों से उस कार्बन पर, जिससे वे आबंधित हैं, चुंबकन स्थानांतरित करता है,
रेडियो-आवृत्ति स्पंदों के एक अनुक्रम द्वारा जिसका अंतिम कोण अनुक्रिया चुनता है; स्पंद यह कैसे करते हैं, इसकी
व्याख्या तृतीय वर्ष के खंड में है। वियुग्मित स्पेक्ट्रम की DEPT-135 और DEPT-90 से तुलना करने पर हर
रेखा C, CH, \ce{CH2} या \ce{CH3} में छँट जाती है।

\begin{omfigure}
\pgfplotsset{dept/.style={width=\linewidth, height=2.9cm, ycomb, x dir=reverse, ymin=-1, ymax=1.15,
  ytick=\empty, axis y line=none, axis x line=middle, x axis line style={-}, tick label style={font=\tiny},
  title style={font=\scriptsize, at={(0.0,0.95)}, anchor=north west}, every axis plot/.append style={thick}}}
\begin{minipage}[t]{0.49\linewidth}\centering
{\small ब्यूटेन-2-ओन}\par
\begin{tikzpicture}
\begin{axis}[dept, xmin=0, xmax=220, xtick={0,50,100,150,200}, title={वियुग्मित}]
\addplot[omDef] table[x=shift, y=dec] {figdata/out/bachelor-2-nmr-spectra-a.dat};
\end{axis}
\end{tikzpicture}\par
\begin{tikzpicture}
\begin{axis}[dept, xmin=0, xmax=220, xtick={0,50,100,150,200}, title={DEPT-135}]
\addplot[omThm] table[x=shift, y=d135] {figdata/out/bachelor-2-nmr-spectra-a.dat};
\end{axis}
\end{tikzpicture}\par
\begin{tikzpicture}
\begin{axis}[dept, xmin=0, xmax=220, xtick={0,50,100,150,200}, title={DEPT-90}]
\addplot[omMeth] table[x=shift, y=d90] {figdata/out/bachelor-2-nmr-spectra-a.dat};
\end{axis}
\end{tikzpicture}
\end{minipage}\hfill
\begin{minipage}[t]{0.49\linewidth}\centering
{\small एथिल बेन्ज़ोएट}\par
\begin{tikzpicture}
\begin{axis}[dept, xmin=0, xmax=220, xtick={0,50,100,150,200}, title={वियुग्मित}]
\addplot[omDef] table[x=shift, y=dec] {figdata/out/bachelor-2-nmr-spectra-b.dat};
\end{axis}
\end{tikzpicture}\par
\begin{tikzpicture}
\begin{axis}[dept, xmin=0, xmax=220, xtick={0,50,100,150,200}, title={DEPT-135}]
\addplot[omThm] table[x=shift, y=d135] {figdata/out/bachelor-2-nmr-spectra-b.dat};
\end{axis}
\end{tikzpicture}\par
\begin{tikzpicture}
\begin{axis}[dept, xmin=0, xmax=220, xtick={0,50,100,150,200}, title={DEPT-90}]
\addplot[omMeth] table[x=shift, y=d90] {figdata/out/bachelor-2-nmr-spectra-b.dat};
\end{axis}
\end{tikzpicture}
\end{minipage}

{\small वियुग्मित \ce{^{13}C} स्पेक्ट्रम (मापे गए विस्थापन और ऊँचाइयाँ, पबकेम आँकड़े), \omterm{def:b2:structure-determination:dept}{DEPT}
अनुक्रियाओं के साथ (चिह्न \cref{prop:b2:structure-determination:dept} से, ऊँचाइयाँ व्यवस्था-रूप)। ब्यूटेन-2-ओन: C=O
209.3 (\omterm{def:b2:structure-determination:dept}{DEPT} में अनुपस्थित), \ce{CH2} 36.9 (नीचे), \ce{CH3} 29.4 और 7.9 (ऊपर)। एथिल बेन्ज़ोएट: C=O 166.5 और
एस्टर धारण करने वाला वलय कार्बन, 130.6, लुप्त हो जाते हैं; तीन \omterm{def:b2:huckel:aromatic}{ऐरोमैटिक} CH (132.8, 129.6, 128.3)
DEPT-90 में रहते हैं; \ce{OCH2} 60.9 नीचे इंगित करता है।}
% figdata: bachelor-2/nmr-spectra.py parts a, b
\end{omfigure}

\section{तकनीकों का संयोजन}

हर तकनीक अपने प्रश्न का उत्तर देती है, और उत्तरों को एक निश्चित क्रम में जोड़ा जाता है: पहले
सूत्र, क्योंकि वह बाक़ी सब को सीमित करता है; फिर क्रियात्मक समूह और ढाँचा; फिर
संयोजन।

\begin{method}[किसी अज्ञात को हल करना]\label{met:b2:structure-determination:solve}
\begin{enumerate}
  \item सूत्र: \omterm{def:b2:mass-spec-atomic:molecular-ion}{आण्विक आयन}, कार्बनों की संख्या के लिए $M+1$, \omterm{def:b2:mass-spec-atomic:isotope-pattern}{समस्थानिक पैटर्न}, \omterm{def:b2:mass-spec-atomic:nitrogen-rule}{नाइट्रोजन नियम}; उपलब्ध हो तो
        यथातथ द्रव्यमान। असंतृप्ति की मात्रा।
  \item अवरक्त: \ce{C=O} (और किस प्रकार का), \ce{O-H}, \ce{N-H}, \ce{C#N}, \qty{3000}{cm^{-1}} से ऊपर \omterm{def:b2:huckel:aromatic}{ऐरोमैटिक} या ऐल्कीन
        \ce{C-H}।
  \item कार्बन-13 और \omterm{def:b2:structure-determination:dept}{DEPT}: कार्बन के प्रकारों की संख्या (सममिति), विस्थापनों से उनके वर्ग, और
        \omterm{def:b2:structure-determination:dept}{DEPT} से C, CH, \ce{CH2}, \ce{CH3}। सूत्र से तुलना कीजिए।
  \item प्रोटॉन NMR: समाकलन, विस्थापन, बहुकताएँ; $n + 1$ नियम से पड़ोसी।
  \item टुकड़ों को जोड़िए; हर संभावित संरचना लिखिए।
  \item हर संभावित संरचना को द्रव्यमान-स्पेक्ट्रम के विखंडों सहित हर आँकड़े से जाँचिए; वह रखिए
        जो उन सबकी व्याख्या करे।
\end{enumerate}
\end{method}

\begin{omfigure}
\begin{tikzpicture}[x=1cm, y=1cm, st/.style={draw=black!70, rounded corners=2pt, fill=white,
  font=\small, align=center, minimum height=0.9cm, inner sep=3pt, text width=2.3cm}]
\node[st, fill=omDef!10] (ms) at (0,0) {MS\\सूत्र, $M+1$};
\node[st] (dbe) at (2.9,0) {असंतृप्ति की\\मात्रा};
\node[st, fill=omDef!10] (ir) at (5.8,0) {IR\\क्रियात्मक समूह};
\node[st, fill=omDef!10] (c13) at (8.7,0) {\ce{^{13}C}, DEPT\\ढाँचा};
\node[st, fill=omDef!10] (h1) at (8.7,-1.8) {\ce{^1H}\\संयोजन};
\node[st] (cand) at (5.8,-1.8) {संभावित\\संरचनाएँ};
\node[st, fill=omThm!10] (chk) at (2.9,-1.8) {हर आँकड़े की\\जाँच};
\node[st] (ans) at (0,-1.8) {संरचना};
\draw[->, thick] (ms) -- (dbe); \draw[->, thick] (dbe) -- (ir); \draw[->, thick] (ir) -- (c13);
\draw[->, thick] (c13) -- (h1); \draw[->, thick] (h1) -- (cand); \draw[->, thick] (cand) -- (chk);
\draw[->, thick] (chk) -- (ans);
\draw[->, thick, densely dashed] (chk.south) -- ++(0,-0.5) -| node[pos=0.25, below, font=\footnotesize]
  {कोई आँकड़ा अव्याख्यात: संभावित संरचनाओं पर वापस} (cand.south);
\end{tikzpicture}
\par\vspace{4pt}

{\small किसी अज्ञात के लिए कार्य का क्रम (\cref{met:b2:structure-determination:solve})। अंतिम चरण
वही है जो सबसे अधिक छोड़ दिया जाता है: जो संरचना एक भी शिखर की व्याख्या न करे, वह अभी उत्तर नहीं है।}
\end{omfigure}

\section{हल किए गए उदाहरण}

\begin{example}[एक कीटोन, C$_5$H$_{10}$O]\label{ex:b2:structure-determination:ketone}
किसी अज्ञात का $M = 86$ है, \qty{1715}{cm^{-1}} के निकट एक IR पट्ट, 208.9, 45.7, 29.8, 17.4 और
13.7~ppm पर कार्बन रेखाएँ, और 2.41 (t, 2H), 2.13 (s, 3H), 1.61 (षट्क, 2H) और 0.91 (t, 3H) पर प्रोटॉन संकेत।
असंतृप्ति की एक कोटि, एक संतृप्त कीटोन (9.7 के निकट कोई ऐल्डिहाइड प्रोटॉन नहीं)। पाँच कार्बन, पाँच रेखाएँ:
कोई सममिति नहीं। 2.13 पर तीन प्रोटॉनों का एकक कार्बोनिल के बगल वाला \ce{CH3} है; त्रिक,
षट्क, त्रिक अनुक्रम एक प्रोपिल समूह है जिसका 2.41 वाला \ce{CH2} कार्बोनिल को छूता है। पेंटेन-2-ओन,
\ce{CH3COCH2CH2CH3}। इसका सममित समावयवी पेंटेन-3-ओन तीन कार्बन रेखाएँ दिखाता और कोई एकक नहीं;
इसके \omterm{def:b2:mass-spec-atomic:spectrum}{द्रव्यमान स्पेक्ट्रम} में 58 पर मैकलैफ़र्टी आयन है जो पेंटेन-3-ओन में नहीं है (\cref{ch:b2:mass-spec-atomic})।
% ledger: c13:pentan-2-one, nmr:pentan-2-one, ir:CO-ketone, ms:pentan-2-one
\end{example}

\begin{example}[एक ऐरोमैटिक एस्टर, C$_9$H$_{10}$O$_2$]\label{ex:b2:structure-determination:ester}
सूत्र \ce{C9H10O2} (असंतृप्ति की पाँच कोटियाँ) वाला एक अज्ञात एक कार्बोनिल पट्ट और सात
कार्बन रेखाएँ दिखाता है: 166.5 (C), 132.8 (CH), 130.6 (C), 129.6 (CH), 128.3 (CH), 60.9 (\ce{CH2}) और 14.3
(\ce{CH3})। इसका प्रोटॉन स्पेक्ट्रम: 8.05 (m, 2H), 7.35--7.63 (m, 3H), 4.37 (q, 2H), 1.38 (t, 3H)। एक
एकल-प्रतिस्थापित बेन्ज़ीन वलय (चार \omterm{def:b2:huckel:aromatic}{ऐरोमैटिक} कार्बन रेखाएँ, जिनमें से दो दो-दो कार्बनों की; पाँच
\omterm{def:b2:huckel:aromatic}{ऐरोमैटिक} प्रोटॉन) चार असंतृप्तियों की व्याख्या करता है, कार्बोनिल पाँचवीं की। चतुष्क--त्रिक युग्म
एक एथिल समूह है; 4.37 पर (और कार्बन में 60.9 पर) उसका \ce{CH2} ऑक्सीजन से आबंधित है। 166.5 पर एस्टर कार्बोनिल
वलय से आबंधित है: उसके बगल वाले 8.05 पर के दो प्रोटॉन उच्चतर विस्थापन की ओर धकेले जाते हैं। एथिल
बेन्ज़ोएट, \ce{C6H5COOCH2CH3}।
% ledger: c13:ethylbenzoate, nmr:ethylbenzoate
\end{example}

\section{जाल}

\begin{itemize}
  \item \emph{विनिमेय प्रोटॉन} (\ce{O-H}, \ce{N-H}) परिवर्ती विस्थापनों पर चौड़े संकेत देते हैं,
        प्रायः अपने पड़ोसियों से युग्मित नहीं होते, और नमूने के साथ \ce{D2O} की एक बूँद हिलाने पर
        लुप्त हो जाते हैं।
  \item \emph{अतिव्यापी संकेत}: बेन्ज़ीन वलय के दो कार्बन, या किसी शृंखला के कई \ce{CH2},
        एक रेखा साझा कर सकते हैं; सममिति पर निष्कर्ष निकालने से पहले रेखाओं को सूत्र से मिलाकर गिनिए।
  \item \emph{अप्रतिबिंबी-स्थानी प्रोटॉन}: किसी त्रिविमजनक केंद्र के बगल वाले \ce{CH2} के दोनों प्रोटॉन
        तुल्य नहीं होते, उनके विस्थापन भिन्न हो सकते हैं और वे आपस में युग्मित होते हैं।
  \item \emph{द्वितीय-कोटि पैटर्न}: जब दो युग्मित प्रोटॉनों के विस्थापन उनके युग्मन स्थिरांक के
        कुछ गुना से अधिक निकट हों, तो बहुक एक-दूसरे की ओर झुकते हैं और $n+1$ नियम विफल होता है;
        उच्च-क्षेत्र स्पेक्ट्रममापी सरल पैटर्न लौटा देता है।
\end{itemize}

\begin{safety}{\ghs[5mm]{GHS02}\,\ghs[5mm]{GHS07}\,\ghs[5mm]{GHS08}}
ब्यूटेन-2-ओन ज्वलनशील है। NMR स्पेक्ट्रममापी का चुंबक सदैव चालू रहता है: उसके पास न इस्पात के औज़ार, न गैस
सिलिंडर, न चुंबकीय कार्ड, और पेसमेकर वाले लोगों का प्रवेश वर्जित।
% ledger: ghs:butanone, ghs:benzylacetate
\end{safety}

\section{अभ्यास}

\begin{exercise}[$\star$]\label{exo:b2:structure-determination:1}
हर ज़ाइलीन (1,2-, 1,3- और
1,4-डाइमेथिलबेन्ज़ीन) का वियुग्मित \ce{^{13}C} स्पेक्ट्रम कितनी रेखाएँ दिखाता है?
\end{exercise}

\begin{exercise}[$\star$]\label{exo:b2:structure-determination:2}
ब्यूटेन-2-ऑल के DEPT-135 और DEPT-90 स्पेक्ट्रमों का पूर्वानुमान कीजिए।
\end{exercise}

\begin{exercise}[$\star$]\label{exo:b2:structure-determination:3}
209, 170, 129, 64 और 14~ppm पर कोई रेखा सबसे संभवतः किस प्रकार के कार्बन की है?
\end{exercise}

\begin{exercise}[$\star$]\label{exo:b2:structure-determination:4}
\ce{C8H8O2} की असंतृप्ति की मात्रा परिकलित कीजिए और बेन्ज़ीन वलय वाली एक संरचना प्रस्तावित कीजिए।
\end{exercise}

\begin{exercise}[$\star\star$]\label{exo:b2:structure-determination:5}
कोई द्रव \ce{C4H8O} \qty{1715}{cm^{-1}} पर एक IR पट्ट और 209, 37, 29 और 8~ppm पर कार्बन रेखाएँ दिखाता है
(DEPT-135: 37 नीचे, 29 और 8 ऊपर, 209 अनुपस्थित)। इसे पहचानिए।
\end{exercise}

\begin{exercise}[$\star\star$]\label{exo:b2:structure-determination:6}
केवल \ce{^{13}C} NMR से आप पेंटेन-2-ओन को पेंटेन-3-ओन से कैसे अलग पहचानेंगे? प्रोटॉन NMR से? \omterm{def:b2:mass-spec-atomic:spectrum}{द्रव्यमान
स्पेक्ट्रममिति} से?
\end{exercise}

\begin{exercise}[$\star\star$]\label{exo:b2:structure-determination:7}
किसी एस्टर का $M = 88$ है और प्रोटॉन संकेत 4.12 (q, 2H), 2.04 (s, 3H) और 1.26 (t, 3H) पर हैं (अभ्यास के
आँकड़े)। इसे पहचानिए।
\end{exercise}

\begin{exercise}[$\star\star$]\label{exo:b2:structure-determination:8}
\ce{^{13}C} NMR द्वारा 1,2-डाइक्लोरोबेन्ज़ीन को 1,4-डाइक्लोरोबेन्ज़ीन से अलग पहचानिए।
\end{exercise}

\begin{exercise}[$\star\star$]\label{exo:b2:structure-determination:9}
2-क्लोरोब्यूटेन में \ce{CH2} समूह के दोनों प्रोटॉन तुल्य क्यों नहीं हैं? इससे
प्रोटॉन स्पेक्ट्रम पर क्या प्रभाव पड़ता है?
\end{exercise}

\begin{exercise}[$\star\star\star$]\label{exo:b2:structure-determination:10}
कोई अज्ञात \ce{C9H10O} \qty{1690}{cm^{-1}} पर एक IR पट्ट, 200.8 (C), 137.0 (C),
132.9 (CH), 128.6 (CH), 128.0 (CH), 31.8 (\ce{CH2}) और 8.3 (\ce{CH3}) पर कार्बन रेखाएँ, और 7.95 (m,
2H), 7.40--7.55 (m, 3H), 2.98 (q, 2H) और 1.22 (t, 3H) पर प्रोटॉन संकेत दिखाता है (अभ्यास के आँकड़े)। इसे पहचानिए, और
निम्न कार्बोनिल तरंग संख्या को समझाइए।
\end{exercise}

\begin{exercise}[$\star\star\star$]\label{exo:b2:structure-determination:11}
चार समावयवी \ce{C5H10O2}: पेंटेनोइक अम्ल, मेथिल ब्यूटेनोएट, एथिल प्रोपेनोएट और प्रोपिल एथेनोएट।
हर एक के लिए उसके IR या प्रोटॉन NMR स्पेक्ट्रम का एक लक्षण दीजिए जो उसे अन्य तीन से अलग करे।
\end{exercise}

\begin{exercise}[$\star\star\star$]\label{exo:b2:structure-determination:12}
एक रिपोर्ट एथिल बेन्ज़ोएट की 166.5~ppm वाली रेखा को एस्टर धारण करने वाले वलय कार्बन को और
130.6~ppm वाली रेखा को कार्बोनिल कार्बन को निर्दिष्ट करती है। \omterm{def:b2:structure-determination:dept}{DEPT} आँकड़ों और विस्थापन चार्ट से दिखाइए कि
यह निर्धारण सही नहीं हो सकता।
\end{exercise}

\section{समस्या: चमेली के चार स्पेक्ट्रम}

\begin{problem}[{सप्ताहांत समस्या --- आण्विक आयन से सूत्र, अवरक्त कार्बोनिल पट्ट,
कार्बन-13 और DEPT, प्रोटॉन स्पेक्ट्रम, और विखंडों के विरुद्ध जाँची गई एक संरचना}]
\label{pb:b2:structure-determination:1}
चमेली की फूलों जैसी गंध वाला एक रंगहीन द्रव निम्नलिखित आँकड़े देता है। \omterm{def:b2:mass-spec-atomic:spectrum}{द्रव्यमान स्पेक्ट्रम} (NIST
वेबबुक): $m/z$ 150 (31.7), 151 (3.1), 108 (100), 91 (71.7), 90 (40.6), 43 (37.6)। अवरक्त (अभ्यास के
आँकड़े): 1740 और \qty{1230}{cm^{-1}} पर प्रबल पट्ट, \qty{3000}{cm^{-1}} से ठीक ऊपर दुर्बल पट्ट,
\qty{3300}{cm^{-1}} के निकट कोई चौड़ा पट्ट नहीं। कार्बन-13, CDCl$_3$ (पबकेम आँकड़े): 170.70, 136.14, 128.56,
128.24, 66.24, 20.82~ppm। प्रोटॉन, CDCl$_3$, \qty{90}{MHz}: 7.33 (m, 5H), 5.09 (s, 2H), 2.06 (s, 3H)।
% ledger: use:benzylacetate, ms:benzylacetate, c13:benzylacetate, nmr:benzylacetate, ir:CO-ester, iso:C12, iso:C13, mass:C12, mass:H1, mass:O16

\begin{center}
\begin{tikzpicture}
\pgfplotsset{dept/.style={width=0.8\linewidth, height=3.2cm, ycomb, x dir=reverse, ymin=0, ymax=1.15,
  xmin=0, xmax=200, ytick=\empty, axis y line=none, axis x line=bottom, x axis line style={-}, tick label style={font=\tiny},
  title style={font=\scriptsize, at={(0.02,0.82)}, anchor=west}, every axis plot/.append style={thick}}}
\begin{axis}[dept, title={वियुग्मित}, xtick={0,20,40,60,80,100,120,140,160,180,200}]
\addplot[omDef] table[x=shift, y=dec] {figdata/out/bachelor-2-nmr-spectra-c.dat};
\end{axis}
\end{tikzpicture}
\end{center}
% figdata: bachelor-2/nmr-spectra.py part c

\medskip
\noindent\textbf{भाग I --- सूत्र।}
\begin{enumerate}
  \item 151/150 अनुपात से कार्बनों की संख्या का अनुमान लगाइए।
  \item \omterm{def:b2:mass-spec-atomic:nitrogen-rule}{नाइट्रोजन नियम} क्या कहता है?
  \item नौ कार्बनों के साथ द्रव्यमान 150 में से क्या बचता है? ऑक्सीजन वाला एक सूत्र प्रस्तावित कीजिए।
  \item \ce{C10H14O}, जिसका नाममात्र द्रव्यमान भी 150 है, कम संभावित क्यों है?
  \item असंतृप्ति की मात्रा परिकलित कीजिए।
  \item चुने गए सूत्र का \omterm{def:b2:mass-spec-atomic:exact-mass}{एकसमस्थानिक द्रव्यमान} परिकलित कीजिए।
\end{enumerate}

\medskip
\noindent\textbf{भाग II --- अवरक्त।}
\begin{enumerate}[resume]
  \item \qty{1740}{cm^{-1}} वाले पट्ट को निर्दिष्ट कीजिए। उसकी स्थिति किस परिवार का संकेत देती है?
  \item \qty{1230}{cm^{-1}} वाले पट्ट को निर्दिष्ट कीजिए।
  \item \qty{3000}{cm^{-1}} से ऊपर के दुर्बल पट्ट क्या दर्शाते हैं?
  \item \qty{3300}{cm^{-1}} के निकट चौड़े पट्ट की अनुपस्थिति क्या अपवर्जित करती है?
  \item क्या कार्बोनिल वलय से संयुग्मित है? उसकी तरंग संख्या से समझाइए।
\end{enumerate}

\medskip
\noindent\textbf{भाग III --- कार्बन-13 और \omterm{def:b2:structure-determination:dept}{DEPT}।}
\begin{enumerate}[resume]
  \item हर रेखा को किसी प्रकार के कार्बन को निर्दिष्ट कीजिए।
  \item DEPT-135 में कौन-सी रेखा नीचे इंगित करती है?
  \item DEPT-135 में कौन-सी रेखाएँ लुप्त हो जाती हैं?
  \item DEPT-90 में कौन-सी रेखाएँ बची रहती हैं?
  \item एकल-प्रतिस्थापित बेन्ज़ीन वलय के कार्बन के कितने प्रकार होते हैं? समझाइए कि केवल दो \omterm{def:b2:huckel:aromatic}{ऐरोमैटिक}
        CH रेखाएँ क्यों दिखती हैं।
  \item 128.24 वाली रेखा सबसे ऊँची है। क्या वह सबसे अधिक कार्बनों को दर्शाती है?
\end{enumerate}

\medskip
\noindent\textbf{भाग IV --- प्रोटॉन NMR और संरचना।}
\begin{enumerate}[resume]
  \item तीनों प्रोटॉन संकेतों को निर्दिष्ट कीजिए।
  \item 5.09 और 2.06 पर संकेत एकक क्यों हैं?
  \item 5.09 पर \ce{CH2} इतने ऊँचे विस्थापन पर क्यों है?
  \item संरचना जोड़िए और उसका नाम दीजिए।
  \item इसके समावयवी मेथिल 2-फ़ेनिलएथेनोएट, \ce{C6H5CH2COOCH3}, का सूत्र भी यही है। कौन-सा आँकड़ा इसे
        अपवर्जित करता है?
  \item 108, 91 और 43 पर विखंडों को समझाइए।
  \item यौगिक का पूर्णतः विभेदित वियुग्मित \ce{^{13}C} स्पेक्ट्रम कितनी रेखाएँ दिखाता, और
        उनमें से कितनी DEPT-135 में नीचे इंगित करतीं?
\end{enumerate}
\end{problem}
